%\section{Magnetic induction field due arbitrarily magnetized sources}
%\label{sec:induction_arbitrary}
%
%Consider a set of magnetic sources enclosed by a volume $v$ in the subsurface.
%Besides, consider that these sources can be represented by the superposition
%of an arbitrary set of $L$ sources with uniform magnetization directions
%(magnetization components), so that the total-magnetization vector
%at any point $(x', y', z')$ within $v$ is given by:
%\begin{equation}
%\mathbf{h}(\mathbf{r}') = \sum\limits_{\ell = 1}^{L} 
%h_{\ell}(\mathbf{r}') \: \hat{\mathbf{u}}(I_{\ell}, D_{\ell}) \quad ,
%\label{eq:h-vector-total}
%\end{equation}
%where $\mathbf{r}' = \left[x' \:\: y' \:\: z' \right]^{\top}$ is the position vector 
%defining a point $(x', y', z')$ inside $v$ and
%$h_{\ell}(\mathbf{r}')$ (in A/m) is the positive magnetization intensity
%of the $\ell$th magnetization component with 
%constant magnetization direction defined by the unit vector
%\begin{equation}
%\hat{\mathbf{u}}(I, D) =
%\begin{bmatrix}
%\hat{u}_{x}(I, D) \\
%\hat{u}_{y}(I, D) \\
%\hat{u}_{z}(I, D)
%\end{bmatrix} =
%\begin{bmatrix}
%\cos(I) \: \cos(D) \\ 
%\cos(I) \: \sin(D) \\
%\sin(I)
%\end{bmatrix} \quad ,
%\label{eq:u-unit-vector}
%\end{equation}
%for the particular case in which the inclination and declination are $I = I_{\ell}$
%and $D = D_{\ell}$.
%This representation of magnetic sources as the superposition of magnetization
%distributions with uniform direction is in perfect agreement with the commonly
%used superposition principle for magnetic fields \cite[e.g., ][ p. 95]{blakely1996}.
%In this case, the magnetic scalar potential produced by the sources at an
%observation point $(x,y,z)$ located above them, outside $v$, is given by
%\cite[e.g., ][ p. 196]{jackson1999}:
%\begin{equation}
%V(\mathbf{r}) = c_{m} \iiint\limits_{v}
%\mathbf{h}(\mathbf{r}')^{\top} \nabla' \frac{1}{\| \mathbf{r} - \mathbf{r}'\|}
%\: dx'dy'dz' \quad ,
%\label{eq:V-volume-integral}
%\end{equation}
%where $\mathbf{r} = \left[x \:\: y \:\: z \right]^{\top}$ is the position vector
%defining an observation point $(x, y, z)$ located outside $v$, 
%$\tfrac{1}{\| \mathbf{r} - \mathbf{r}'' \|}$ is the inverse distance function, 
%$c_{m} = 10^{9} \frac{\mu_{0}}{4\pi}$ (H/m), $\mu_{0}$ is the
%magnetic permeability at vacuum, $\| \cdot \|$ denotes the Euclidean norm and $\nabla'$ is the gradient
%operator with respect to the coordinates of $(x', y', z')$.
%The resultant magnetic scalar potential $V(\mathbf{r})$ 
%(equation \ref{eq:V-volume-integral}) can be rewritten in terms of the
%magnetic scalar potentials $V_{\ell}(\mathbf{r})$ produced by each magnetization 
%component as follows:
%\begin{equation}
%V(\mathbf{r}) = \sum\limits_{\ell = 1}^{L}
%V_{\ell}(\mathbf{r}) \quad ,
%\label{eq:V-summation}
%\end{equation}
%where
%\begin{equation}
%V_{\ell}(\mathbf{r}) = -\nabla \Theta_{\ell}(\mathbf{r})^{\top} 
%\hat{\mathbf{u}}(I_{\ell}, D_{\ell}) \quad ,
%\label{eq:V-l}
%\end{equation}
%$\nabla$ is the gradient operator with respect to the coordinates of $(x,y,z)$ and
%$\Theta_{\ell}(\mathbf{r})$ is an harmonic function given by
%\begin{equation}
%\Theta_{\ell}(\mathbf{r}) = c_{m} \iiint\limits_{v_{\ell}}
%h_{\ell}(\mathbf{r}') \: \frac{1}{\| \mathbf{r} - \mathbf{r}'\|}
%\: dx'dy'dz' \quad .
%\label{eq:Theta-l}
%\end{equation}
%Note that the magnetic scalar potential $V_{\ell}(\mathbf{r})$ (equation \ref{eq:V-l}) 
%is the negative directional derivative of the harmonic function $\Theta_{\ell}(\mathbf{r})$
%(equation \ref{eq:Theta-l}) along the uniform magnetization direction
%$\hat{\mathbf{u}}(I_{\ell}, D_{\ell})$ (equation \ref{eq:u-unit-vector})
%of the $\ell$th magnetization component. 
%
%By computing the negative gradient of $V(\mathbf{r})$ 
%(equation \ref{eq:V-volume-integral}), we obtain the resultant 
%magnetic induction field $\mathbf{B}(\mathbf{r})$ produced by the 
%sources at the observation point $(x,y,z)$:
%\begin{equation}
%\mathbf{B}(\mathbf{r}) = -\nabla V(\mathbf{r}) = \sum\limits_{\ell = 1}^{L}
%\mathbf{B}_{\ell}(\mathbf{r}) \quad ,
%\label{eq:B-summation}
%\end{equation}
%where 
%\begin{equation}
%\mathbf{B}_{\ell}(\mathbf{r}) = \mathbf{H}_{\ell}(\mathbf{r}) 
%\: \hat{\mathbf{u}}(I_{\ell}, D_{\ell})
%\label{eq:B-l}
%\end{equation}
%is the magnetic induction field produced by the $\ell$th magnetization component
%and 
%\begin{equation}
%\mathbf{H}_{\ell}(\mathbf{r}) = 
%\begin{bmatrix}
%\partial_{xx} \Theta_{\ell}(\mathbf{r}) & \partial_{xy} \Theta_{\ell}(\mathbf{r}) & \partial_{xz} \Theta_{\ell}(\mathbf{r}) \\
%\partial_{xy} \Theta_{\ell}(\mathbf{r}) & \partial_{yy} \Theta_{\ell}(\mathbf{r}) & \partial_{yz} \Theta_{\ell}(\mathbf{r}) \\
%\partial_{xz} \Theta_{\ell}(\mathbf{r}) & \partial_{yz} \Theta_{\ell}(\mathbf{r}) & \partial_{zz} \Theta_{\ell}(\mathbf{r})
%\end{bmatrix}
%\label{eq:H-l}
%\end{equation}
%is the Hessian matrix of $\Theta_{\ell}(\mathbf{r})$ (equation \ref{eq:Theta-l}), 
%with $\partial_{\alpha\beta} \Theta_{\ell}(\mathbf{r}) \equiv 
%\frac{\partial^{2}}{\partial\alpha \, \partial\beta} \Theta_{\ell}(\mathbf{r})$,
%$\alpha = x, y, z$, $\beta = x, y, z$.
%
%From equations \ref{eq:B-summation} and \ref{eq:B-l}, it is easy to show that
%the component of the magnetic induction field along an arbitrary direction $\hat{\mathbf{u}}(I_{t}, D_{t})$ 
%(equation \ref{eq:u-unit-vector}) with constant inclination $I = I_{t}$ and 
%declination $D = D_{t}$ is given by:
%\begin{equation}
%\Delta\tilde{T}_{\ell}(\mathbf{r}) 
%= \hat{\mathbf{u}}(I_{t}, D_{t})^{\top} \mathbf{B}_{\ell}(\mathbf{r}) \quad ,
%\label{eq:Delta-T-l}
%\end{equation}
%if we consider only the field $\mathbf{B}_{\ell}(\mathbf{r})$ produced by the $\ell$th 
%magnetization component, or
%\begin{equation}
%\Delta\tilde{T}(\mathbf{r}) 
%= \hat{\mathbf{u}}(I_{t}, D_{t})^{\top} \mathbf{B}(\mathbf{r}) \quad ,
%\label{eq:Delta-T}
%\end{equation}
%if we consider the resultant field $\mathbf{B}(\mathbf{r})$. 
%The field components $\Delta\tilde{T}_{\ell}(\mathbf{r})$ (equation \ref{eq:Delta-T-l})
%and $\Delta\tilde{T}(\mathbf{r})$ (equation \ref{eq:Delta-T})
%are harmonic functions commonly used in applied geophysics to approximate the total-field anomaly 
%produced by magnetized sources having uniform and nonuniform magnetization directions,
%respectively, where $\hat{\mathbf{u}}(I_{t}, D_{t})$ defines the constant direction of the main geomagnetic 
%field in the study area \cite[][ p. 179]{blakely1996}.

\section{Magnetic induction field due sources with uniform magnetization direction}
\label{sec:induction_uniform_direction}


Consider a set of magnetic sources enclosed by a finite volume $v$ located 
in the subsurface, referred to a topocentric Cartesian coordinate system with axes 
$x$, $y$, and $z$ pointing to North, East and down, respectively. 
Consider also that the total-magnetization vector at each point $(x', y', z') \in v$ 
within the sources, with an associated position vector 
$\mathbf{r}' = \left[x' \:\: y' \:\: z' \right]^{\top}$, is given by:
\begin{equation}
\mathbf{h}(\mathbf{r}') = h(\mathbf{r}') \, \hat{\mathbf{u}}(I_{h}, D_{h}) \: ,
\label{eq:total-magnetization-vector}
\end{equation}
where $h(\mathbf{r}')$ is a positive scalar function representing a spatially-varying 
total-magnetization intensity (in A/m) and $\hat{\mathbf{u}}(I_{h}, D_{h})$ 
is a unit vector defining a uniform total-magnetization direction according to
the vector function
\begin{equation}
\hat{\mathbf{u}}(I, D) =
\begin{bmatrix}
\hat{u}_{x}(I, D) \\
\hat{u}_{y}(I, D) \\
\hat{u}_{z}(I, D)
\end{bmatrix} =
\begin{bmatrix}
\cos(I) \: \cos(D) \\ 
\cos(I) \: \sin(D) \\
\sin(I)
\end{bmatrix} \quad ,
\label{eq:u-unit-vector}
\end{equation}
with constant inclination $I = I_{h}$ and declination $D = D_{h}$ values.
In this case, the magnetic scalar potential produced by the set of sources at an
observation point $(x,y,z)$ located above them, outside $v$, is given by
\cite[e.g., ][ p. 196]{jackson1999}:
\begin{equation}
V(\mathbf{r}) 
= -\nabla \Theta(\mathbf{r})^{\top} 
\hat{\mathbf{u}}(I_{h}, D_{h}) \: ,
\label{eq:scalar-potential-V}
\end{equation}
which represents a negative directional derivative of the harmonic function
\begin{equation}
\Theta(\mathbf{r}) = c_{m} \iiint\limits_{v}
h(\mathbf{r}') \: \frac{1}{\| \mathbf{r} - \mathbf{r}'\|}
\: dx'dy'dz' \: ,
\label{eq:scalar-potential-Theta}
\end{equation}
where the superscript $^{\top}$ stands for transposition; 
$\mathbf{r} = \left[x \:\: y \:\: z \right]^{\top}$ is the position vector 
associated with the observation point $(x, y, z)$; 
$\nabla$ is the gradient operator with respect to the coordinates of $\mathbf{r}$;
$\tfrac{1}{\| \mathbf{r} - \mathbf{r}'' \|}$ is the inverse distance function; 
$\| \cdot \|$ denotes the Euclidean norm; $c_{m} = 10^{9} \frac{\mu_{0}}{4\pi}$ (in H/m);
and $\mu_{0}$ is the magnetic permeability at vacuum. 
It is important noting that the harmonic function $\Theta(\mathbf{r})$ (equation
\ref{eq:scalar-potential-Theta}) represents the gravitational potential that would be
produced by the sources if they had a positive density distribution proportional
to the total-magnetization intensity $h(\mathbf{r}')$ 
(equation \ref{eq:total-magnetization-vector}). Because of that,
$\Theta(\mathbf{r})$ is strictly positive at all observation points 
$\mathbf{r}$ located outside the volume $v$ of the sources.

The magnetic induction field produced by the sources at the observation 
point $\mathbf{r}$ is defined by the negative gradient of the scalar potential
$V(\mathbf{r})$ (equation \ref{eq:scalar-potential-V}) as follows:
\begin{equation}
\mathbf{B}(\mathbf{r}) 
= \mathbf{H}_{\Theta}(\mathbf{r}) \: \hat{\mathbf{u}}(I_{h}, D_{h}) \: ,
\label{eq:B-field}
\end{equation}
where 
\begin{equation}
\mathbf{H}_{\Theta}(\mathbf{r}) = 
\begin{bmatrix}
\partial_{xx} \Theta(\mathbf{r}) & 
\partial_{xy} \Theta(\mathbf{r}) & 
\partial_{xz} \Theta(\mathbf{r}) \\
\partial_{xy} \Theta(\mathbf{r}) & 
\partial_{yy} \Theta(\mathbf{r}) & 
\partial_{yz} \Theta(\mathbf{r}) \\
\partial_{xz} \Theta(\mathbf{r}) & 
\partial_{yz} \Theta(\mathbf{r}) & 
\partial_{zz} \Theta(\mathbf{r})
\end{bmatrix}
\label{eq:H-Theta}
\end{equation}
is the Hessian matrix of $\Theta(\mathbf{r})$ (equation \ref{eq:scalar-potential-Theta}), 
with elements defined by the second derivatives 
$\partial_{\alpha\beta} \Theta(\mathbf{r}) \equiv 
\frac{\partial^{2}}{\partial\alpha \, \partial\beta} \Theta(\mathbf{r})$,
$\alpha = x, y, z$, $\beta = x, y, z$. 
It is easy to verify from equation \ref{eq:B-field} that the $\alpha$-component,
$\alpha = x, y, z$, of the magnetic induction field is given by:
\begin{equation}
B_{\alpha}(\mathbf{r}) 
= \partial_{\alpha} \nabla \Theta(\mathbf{r})^{\top} \hat{\mathbf{u}}(I_{h}, D_{h}) \: ,
\:\: \alpha = x, y, z \: ,
\label{eq:B-alpha}
\end{equation}
where $\partial_{\alpha} \nabla \Theta(\mathbf{r})^{\top}$ denotes a row of the
Hessian matrix $\mathbf{H}_{\Theta}(\mathbf{r})$ (equation \ref{eq:H-Theta}).

Finally, there are two important quantities for geophysical applications.
The first is the component of the magnetic induction field $\mathbf{B}(\mathbf{r})$
(equation \ref{eq:B-field}) along an arbitrary direction 
$\hat{\mathbf{u}}(I_{t}, D_{t})$ (equation \ref{eq:u-unit-vector}) with constant 
inclination $I = I_{t}$ and declination $D = D_{t}$:
\begin{equation}
B_{t}(\mathbf{r}) 
= \hat{\mathbf{u}}(I_{t}, D_{t})^{\top} \mathbf{B}(\mathbf{r}) \: ,
\label{eq:B-t}
\end{equation}
which is a harmonic function commonly used in applied geophysics to approximate the
total-field anomaly produced by magnetized sources having uniform magnetization direction,
where $\hat{\mathbf{u}}(I_{t}, D_{t})$ defines the presumed constant direction of the main
geomagnetic field at the study area \cite[][ p. 179]{blakely1996}. Because of that,
we conveniently call $B_{t}(\mathbf{r})$ (equation \ref{eq:B-t}) 
\textit{approximated total-field anomaly}.
The second important quantity for geophysical applications is the vertical component 
of the magnetic induction field $\mathbf{B}(\mathbf{r})$ (equation \ref{eq:B-field})
produced by the sources for the particular case in which they have a vertical 
magnetization direction:
\begin{equation}
B_{r}(\mathbf{r}) = \partial_{zz} \Theta(\mathbf{r}) \: ,
\label{eq:B-r}
\end{equation}
where $\partial_{zz} \Theta(\mathbf{r})$ is the second derivative of the
harmonic function $\Theta(\mathbf{r})$ (equation \ref{eq:scalar-potential-Theta}) with
respect to the coordinate $z$ of the observation point.
Equation \ref{eq:B-r} is deduced from equations \ref{eq:B-field} and 
\ref{eq:B-t} for the particular case in which $I_{t} = I_{h} = 90^{\circ}$ and
is commonly used to represent the magnetic induction field 
that would be produced by the sources if they had a purely induced and vertical
magnetization. 
Differently from generic total-field anomalies 
\cite[e.g.,][ p. 178]{blakely1996}, which have a ``skewed'' pattern 
oscillating from negative to positive values and do not have their maxima placed right 
above the magnetic sources, equation \ref{eq:B-r} defines a simpler anomaly that is
better for interpretation because it is predominantly positive right above the sources and
decays to zero towards their boundaries. \citet{baranov1957} was the first
to propose a method, called \textit{reduction to the pole} (RTP), to transform generic
total-field anomalies into $B_{r}(\mathbf{r})$ (equation \ref{eq:B-r}),
conveniently called here \textit{RTP anomaly}.


\section{Generic planar equivalent layer of dipoles}
\label{eq:generic_planar_layer_dipoles}


Consider a generic planar layer of dipoles located above the magnetic sources, at 
$z = z_{c}$,
with uniform magnetization direction $\hat{\mathbf{u}}(I_{p}, D_{p})$ 
(equation \ref{eq:u-unit-vector}) defined by the constant inclination $I_{p}$ and
declination $D_{p}$.
Now, consider that this planar layer exactly reproduces the approximated total-field anomaly
$B_{t}(\mathbf{r})$ (equation \ref{eq:B-t}) at observation points $\mathbf{r}$ 
located above the layer, with $z < z_{c}$. In this case, we may write the following 
integral equation:
\begin{equation}
B_{t}(\mathbf{r}) = 
\int\limits_{-\infty}^{\infty} \int\limits_{-\infty}^{\infty}
p(\mathbf{r}'') \:\: 
\partial_{tp} \Psi(\mathbf{r}, \mathbf{r}'') \: 
dx''dy'' \: , \:\: z < z_{c} \: ,
\label{eq:B-t-generic-layer}
\end{equation}
where $\mathbf{r}'' = \left[x'' \:\: y'' \:\: z_{c} \right]^{\top}$ is the 
position vector defining a point $(x'', y'', z_{c})$ located on the plane 
$z = z_{c}$; $p(\mathbf{r}'')$ is an unknown scalar function representing the 
magnetic moment distribution on the planar layer;
\begin{equation}
\partial_{tp} \Psi(\mathbf{r}, \mathbf{r}'') \equiv
\hat{\mathbf{u}}(I_{t}, D_{t})^{\top} 
\mathbf{H}_{\Psi}(\mathbf{r}, \mathbf{r}'') \: 
\hat{\mathbf{u}}(I_{p}, D_{p}) 
\label{eq:Dtp-Psi}
\end{equation}
denotes the second directional derivative of the inverse distance function
\begin{equation}
\Psi(\mathbf{r}, \mathbf{r}'') = 
\frac{1}{2\pi} \: 
\frac{1}{\| \mathbf{r} - \mathbf{r}'' \|}
\label{eq:Psi}
\end{equation}
along the uniform magnetization direction $\hat{\mathbf{u}}(I_{p}, D_{p})$ of the planar
layer and the presumed constant direction $\hat{\mathbf{u}}(I_{t}, D_{t})$ 
of the main geomagnetic field at the study area; and
\begin{equation}
\mathbf{H}_{\Psi}(\mathbf{r}, \mathbf{r}'') = 
\begin{bmatrix}
\partial_{xx} \Psi(\mathbf{r}, \mathbf{r}'') & 
\partial_{xy} \Psi(\mathbf{r}, \mathbf{r}'') & 
\partial_{xz} \Psi(\mathbf{r}, \mathbf{r}'') \\
\partial_{xy} \Psi(\mathbf{r}, \mathbf{r}'') & 
\partial_{yy} \Psi(\mathbf{r}, \mathbf{r}'') & 
\partial_{yz} \Psi(\mathbf{r}, \mathbf{r}'') \\
\partial_{xz} \Psi(\mathbf{r}, \mathbf{r}'') & 
\partial_{yz} \Psi(\mathbf{r}, \mathbf{r}'') & 
\partial_{zz} \Psi(\mathbf{r}, \mathbf{r}'')
\end{bmatrix}
\label{eq:H-Psi}
\end{equation}
is the Hessian matrix of the inverse distance function 
$\Psi(\mathbf{r}, \mathbf{r}'')$ (equation \ref{eq:Psi}), with elements defined
by second derivatives $\partial_{\alpha\beta} \Psi(\mathbf{r}, \mathbf{r}'') \equiv 
\frac{\partial^{2}}{\partial\alpha \, \partial\beta} \Psi(\mathbf{r}, \mathbf{r}'')$,
$\alpha = x, y, z$, $\beta = x, y, z$, with respect to the coordinates of
the observation point $\mathbf{r}$.

Methods based on equivalent layer for dealing with total-field anomaly data 
have been traditionally developed in space domain 
\cite[e.g.,][]{silva_1986, oliveirajr-etal_2013}
by numerically solving the integral equation \ref{eq:B-t-generic-layer} 
for the unknown physical-property distribution $p(\mathbf{r}'')$, 
with predefined constant values for the inclination $I = I_{p}$ and declination $D = D_{p}$. 
Those methods usually assume that total-field anomaly data 
can be satisfactorily approximated by the harmonic function $B_{t}(\mathbf{r})$
(equation \ref{eq:B-t-generic-layer}).
After obtaining a numerical solution for the physical-property distribution 
$p(\mathbf{r}'')$, it is used to compute a potential-field transformation. 
If the desired transformation does not
cause phase changes in the magnetic data, such as the upward/downward continuation, 
then the predefined constant values for the inclination $I_{p}$ and declination $D_{p}$ 
may be arbitrarily chosen. 
Otherwise, the $I_{p}$ and $D_{p}$ must assume the same constant inclination and declination
values defining the total-magnetization direction of the true sources. An example of 
transformation that produces phase changes in the magnetic data is the RTP.

What follows aims at deducing solutions for the integral equation \ref{eq:B-t-generic-layer}
in space and Fourier domains and proving the existence of a planar equivalent 
layer of dipoles that is located above the sources and exactly reproduces the magnetic 
scalar potential $V(\mathbf{r})$ (equation \ref{eq:scalar-potential-V}), magnetic induction 
field $\mathbf{B}(\mathbf{r})$ (equation \ref{eq:B-field}),
approximated total-field anomaly $\mathbf{B}_{t}(\mathbf{r})$ (equation \ref{eq:B-t}) 
and RTP anomaly $\mathbf{B}_{r}(\mathbf{r})$ (equation \ref{eq:B-r})
yielded by sources with uniform total-magnetization direction. 


\section{Planar equivalent layer of dipoles}
\label{sec:eqlayer}


\subsection{Continuous analytical solutions in space domain with true magnetization direction}
\label{subsec:space_domain_true_mag}


The harmonic function $\Theta(\mathbf{r})$ (equation \ref{eq:scalar-potential-Theta}) 
can be alternatively defined in terms of its vertical 
derivative on a plane $z = z_{c}$ located between the observation points and the 
magnetic sources as follows \cite[e.g.,][]{roy1962, reis_etal2020}:
\begin{equation}
\Theta(\mathbf{r}) =  
\int\limits_{-\infty}^{\infty} \int\limits_{-\infty}^{\infty}
\partial_{z} \Theta(\mathbf{r}'') \:\: 
\Psi(\mathbf{r}, \mathbf{r}'') \: dx''dy'' \: , \quad z < z_{c} \: ,
\label{eq:Theta-Neumann}
\end{equation}
where $\partial_{z} \Theta(\mathbf{r}'') \equiv 
\frac{\partial}{\partial z} \Theta(\mathbf{r}'')$ denotes the partial
derivative, evaluated at $(x'',y'',z_{c})$, of the harmonic function
$\Theta(\mathbf{r})$ (equation \ref{eq:scalar-potential-Theta}) with respect 
to the coordinate $z$ and $\Psi(\mathbf{r}, \mathbf{r}'')$ is the inverse 
distance function defined by equation \ref{eq:Psi}.
Equation \ref{eq:Theta-Neumann} represents the solution of the 
\textit{Neumann’s problem} or the \textit{second boundary value problem of potential
theory} \cite[][ p. 246]{kellogg1967} on a plane.
It is also important noting that equation \ref{eq:Theta-Neumann} represents a
two-dimensional convolution of the functions $\partial_{z} \Theta(\mathbf{r}'')$
and $\Psi(\mathbf{r}, \mathbf{r}'')$ (equation \ref{eq:Psi}) on the
plane $z = z_{c}$.

By computing the negative directional derivative of $\Theta(\mathbf{r})$ 
(equation \ref{eq:Theta-Neumann}) along the uniform magnetization direction
$\hat{\mathbf{u}}(I_{h}, D_{h})$ (equation \ref{eq:total-magnetization-vector}) 
of the sources, we obtain the following alternative expression for the magnetic 
scalar potential $V(\mathbf{r})$ (equation \ref{eq:scalar-potential-V}):
\begin{equation}
V(\mathbf{r}) = -
\int\limits_{-\infty}^{\infty} \int\limits_{-\infty}^{\infty}
\partial_{z} \Theta(\mathbf{r}'') \:\:
\partial_{h} \Psi(\mathbf{r}, \mathbf{r}'') \: dx''dy'' \: , \quad z < z_{c} \: ,
\label{eq:V-Neumann}
\end{equation}
where
\begin{equation}
\partial_{h} \Psi(\mathbf{r}, \mathbf{r}'') \equiv 
\nabla \Psi(\mathbf{r}, \mathbf{r}'')^{\top}
\hat{\mathbf{u}}(I_{h}, D_{h})
\label{eq:Dh-Psi}
\end{equation}
represents the directional derivative of the inverse distance function 
$\Psi(\mathbf{r}, \mathbf{r}'')$ (equation \ref{eq:Psi}) along the uniform
total-magnetization direction $\hat{\mathbf{u}}(I_{h}, D_{h})$ of the true sources
(equation \ref{eq:total-magnetization-vector}).
From the physical point of view, equations \ref{eq:V-Neumann} and \ref{eq:Dh-Psi} 
define a planar layer of dipoles having the same uniform total-magnetization direction 
$\hat{\mathbf{u}}(I_{h}, D_{h})$ of the true sources (equation \ref{eq:total-magnetization-vector}) 
and a positive magnetic-moment intensity distribution 
$\partial_{z} \Theta(\mathbf{r}'')$, which is proportional to the 
pseudogravity anomaly produced by the true sources at the plane $z = z_{c}$. 
\citet{reis_etal2020} have followed a similar approach to deduce this 
all-positive magnetic-moment distribution for interpreting 
total-field anomaly data.
Note that the planar layer of dipoles exactly reproduces the magnetic scalar potential
$V(\mathbf{r})$ (equation \ref{eq:scalar-potential-V}) yielded by the true sources at 
all observation points $\mathbf{r}$ with $z < z_{c}$. 
By computing the negative gradient of the magnetic scalar potential $V(\mathbf{r})$ 
(equation \ref{eq:V-Neumann}), we can also show that the 
magnetic induction field yielded by the planar layer of dipoles 
is given by:
\begin{equation}
\mathbf{B}(\mathbf{r}) =
\int\limits_{-\infty}^{\infty} \int\limits_{-\infty}^{\infty}
\partial_{z} \Theta(\mathbf{r}'') \:\: 
\nabla \partial_{h} \Psi(\mathbf{r}, \mathbf{r}'') \: dx''dy'' \: , \:\: z < z_{c} \: ,
\label{eq:B-field-Neumann}
\end{equation}
where
\begin{equation}
\nabla \partial_{h} \Psi(\mathbf{r}, \mathbf{r}'') \equiv
\mathbf{H}_{\Psi}(\mathbf{r}, \mathbf{r}'') \: 
\hat{\mathbf{u}}(I_{h}, D_{h}) \: .
\label{eq:nabla-Dh-Psi}
\end{equation}
The $\alpha$-component, $\alpha = x, y, z$, of the magnetic induction field
can then be defined as follows:
\begin{equation}
B_{\alpha}(\mathbf{r}) = 
\int\limits_{-\infty}^{\infty} \int\limits_{-\infty}^{\infty}
\partial_{z} \Theta(\mathbf{r}'') \:\: 
\partial_{\alpha h} \Psi(\mathbf{r}, \mathbf{r}'') \: 
dx''dy'' \: , \:\: z < z_{c} \: , \:\: \alpha = x, y, z \: ,
\label{eq:B-alpha-Neumann}
\end{equation}
where $\partial_{\alpha h} \Psi(\mathbf{r}, \mathbf{r}'')$ denotes the second
derivative
\begin{equation}
\partial_{\alpha h} \Psi(\mathbf{r}, \mathbf{r}'') \equiv
\partial_{\alpha} \nabla \Psi (\mathbf{r}, \mathbf{r}'')^{\top} \:
\hat{\mathbf{u}}(I_{h}, D_{h}) \: , \:\: \alpha = x, y, z \: ,
\label{eq:Dalphah-Psi}
\end{equation}
with $\partial_{\alpha} \nabla \Psi (\mathbf{r}, \mathbf{r}'')^{\top}$ representing
a row of the Hessian matrix $\mathbf{H}_{\Psi}$ (equation \ref{eq:H-Psi}).
It is also relatively easy to verify from equations 
\ref{eq:B-t} and \ref{eq:B-field-Neumann} that the approximated total-field 
anomaly produced by the planar layer of dipoles is given by:
\begin{equation}
B_{t}(\mathbf{r}) = 
\int\limits_{-\infty}^{\infty} \int\limits_{-\infty}^{\infty}
\partial_{z} \Theta(\mathbf{r}'') \:\: 
\partial_{th} \Psi(\mathbf{r}, \mathbf{r}'') \: 
dx''dy'' \: , \:\: z < z_{c} \: ,
\label{eq:B-t-Neumann}
\end{equation}
where
\begin{equation}
\partial_{th} \Psi(\mathbf{r}, \mathbf{r}'') \equiv
\hat{\mathbf{u}}(I_{t}, D_{t})^{\top} 
\mathbf{H}_{\Psi}(\mathbf{r}, \mathbf{r}'') \: 
\hat{\mathbf{u}}(I_{h}, D_{h}) 
\label{eq:Dth-Psi}
\end{equation}
denotes the second directional derivative of the inverse distance function
$\Psi(\mathbf{r}, \mathbf{r}'')$ (equation \ref{eq:Psi}) 
along the directions defined by $\hat{\mathbf{u}}(I_{h}, D_{h})$ 
(equation \ref{eq:total-magnetization-vector}) and 
$\hat{\mathbf{u}}(I_{t}, D_{t})$ (equation \ref{eq:B-t}).

Equations \ref{eq:V-Neumann}--\ref{eq:Dth-Psi} show that not only the magnetic scalar
potential $V(\mathbf{r})$ (equation \ref{eq:scalar-potential-V}), but also the magnetic
induction field $\mathbf{B}(\mathbf{r})$ (equation \ref{eq:B-field}) and 
approximated total-field anomaly $B_{t}(\mathbf{r})$ (equation \ref{eq:B-t}) 
yielded by the true sources are exactly reproduced by the same planar layer 
of dipoles located above the sources, at $z = z_{c}$, with uniform magnetization 
direction $\hat{\mathbf{u}}(I_{h}, D_{h})$ and positive magnetic-moment intensity
distribution $\partial_{z} \Theta(\mathbf{r}'')$. 
Note that $p(\mathbf{r}'') = \partial_{z} \Theta(\mathbf{r}'')$ represents
the analytical solution of the integral equation \ref{eq:B-t-generic-layer} for
the particular case in which the equivalent layer has the same uniform magnetization direction $\hat{\mathbf{u}}(I_{h}, D_{h})$ (equation \ref{eq:total-magnetization-vector}) 
of the true sources, i.e., for the case in which $I_{p} = I_{h}$ and $D_{p} = D_{h}$. 
Equations \ref{eq:V-Neumann}--\ref{eq:Dth-Psi} are significant in at least three 
major aspects because they show that:
(i) the planar equivalent layer of dipoles is not a purely 
mathematical expression for magnetic-data processing, but a solution of a
boundary value problem of potential theory that is intrinsically related to
the magnetic field produced by sources with uniform magnetization direction;
(ii) the uniform magnetization direction of the planar equivalent layer is 
equal to that of the true sources, regardless of whether 
this direction is normal to the planar layer or not;
(iii) the magnetic-moment distribution on the planar layer 
is strictly positive at all points $(x'',y'',z_{c})$ and is proportional to the
pseudogravity anomaly produced by the sources.

Finally, we may also verify from equations \ref{eq:B-r} and \ref{eq:Theta-Neumann} 
that the RTP anomaly produced by the planar equivalent layer is given by:
\begin{equation}
B_{r}(\mathbf{r}) = 
\int\limits_{-\infty}^{\infty} \int\limits_{-\infty}^{\infty}
\partial_{z} \Theta(\mathbf{r}'') \:\: 
\partial_{zz} \Psi(\mathbf{r}, \mathbf{r}'') \: 
dx''dy'' \: , \:\: z < z_{c} \: .
\label{eq:B-r-Neumann}
\end{equation}
It is important to stress that equation \ref{eq:B-r-Neumann} represents the 
theoretical basis for computing the RTP in space domain, from total-field
anomaly data, via equivalent-layer technique.
This is a two-step approach that consists in first considering that the observed 
total-field anomaly data can be satisfactorily represented by 
$B_{t}(\mathbf{r})$ and then numerically solving the integral 
equation \ref{eq:B-t-Neumann} to obtain a discrete version of the positive physical 
property distribution $p(\mathbf{r}'') = \partial_{z} \Theta(\mathbf{r}'')$.
In the second step, the estimated $p(\mathbf{r}'')$ is used in 
a discrete version of equation \ref{eq:B-r-Neumann} to compute the RTP anomaly 
$B_{r}(\mathbf{r})$. It is important noting that equations \ref{eq:B-t-Neumann}
and \ref{eq:B-r-Neumann} show, for the first time, the main 
theoretical assumptions supporting the space-domain RTP: 
(i) There exists a planar equivalent layer of dipoles that is located at a coordinate 
$z = z_{c}$ above the sources, has the same uniform magnetization direction of the sources
and approximates the observed total-field anomaly at observation points;
(ii) The same planar equivalent layer of dipoles reproducing the observed total-field
anomaly data is able to reproduce the RTP anomaly $B_{r}(\mathbf{r})$.
These two theoretical assumptions have being implicitly used, but not clearly stated 
neither in the pioneer method proposed by \cite{silva_1986}, nor in the succeeding 
methods \cite[e.g.,][]{leao-silva_1989, oliveirajr-etal_2013, li_etal_2014, reis_etal2020}
proposed to compute the RTP in space-domain, via equivalent-layer technique.


\subsection{Continuous analytical solutions in Fourier domain with true magnetization direction}
\label{subsec:Fourier_domain_true_mag}


Several methods for magnetic-data processing have been developed in the Fourier domain 
for the particular case in which the magnetic data rely on a
planar and regular grid at constant $z = z_{0}$ 
\citep[][ ch 11 and 12]{blakely1996}. In this case, they 
can take advantage of the well-known computational efficiency of 
Fast Fourier Transform (FFT) algorithms.
Most of these methods is based on computing the two-dimensional Fourier transform
(2DFT) of the approximated total-field anomaly $B_{t}(\mathbf{r})$ (equation \ref{eq:B-t})
with respect to the horizontal coordinates $x$ and $y$ of the observation points, according
to \citet{gunn1975} or \citet[][ p. 298]{blakely1996}, for example. 
It is important to stress that equation \ref{eq:B-t} is defined in terms of volume 
integrals that depend on the position, shape and total-magnetization distribution of the
true sources. 
The main goal here is presenting the 2DFT of equations \ref{eq:B-alpha-Neumann},
\ref{eq:B-t-Neumann} and \ref{eq:B-r-Neumann} for the case in which the magnetic data
are on the plane $z = z_{0}$ and show that they are equivalent to the 2DFT of 
the corresponding equations defined by volume integrals (equations 
\ref{eq:B-alpha}--\ref{eq:B-r}).
To do this, we follow the same approach presented by \citet[][ ch. 11]{blakely1996}.

Let us first consider that the 2D convolutions representing the 
$\alpha$-component $B_{\alpha}(\mathbf{r})$ (equation \ref{eq:B-alpha-Neumann}),
approximated total-field anomaly $B_{t}(\mathbf{r})$ (equation \ref{eq:B-t-Neumann}) 
and RTP anomaly $B_{r}(\mathbf{r})$ (equation \ref{eq:B-r-Neumann}) are 
evaluated at observation points $\mathbf{r}_{0} = \left[x \:\: y \:\: z_{0} \right]^{\top}$ 
located on a horizontal plane with constant vertical coordinate $z = z_{0}$. 
Then, we use the Fourier-Convolution Theorem \cite[][ p. 268]{blakely1996} to obtain:
\begin{equation}
\mathcal{F}\left[ B_{\alpha}(\mathbf{r}_{0}) \right] = 
\mathcal{F}\left[ \partial_{z} \Theta(\mathbf{r}'') \right] \: 
\mathcal{F}\left[ \partial_{\alpha h} \Psi(\mathbf{r}_{0}, \mathbf{r}'') \right] \: ,
\:\: \alpha = x, y, z \: ,
\label{eq:Fourier-B-alpha}
\end{equation}
\begin{equation}
\mathcal{F}\left[ B_{t}(\mathbf{r}_{0}) \right] = 
\mathcal{F}\left[ \partial_{z} \Theta(\mathbf{r}'') \right] \: 
\mathcal{F}\left[ \partial_{th} \Psi(\mathbf{r}_{0}, \mathbf{r}'') \right]
\label{eq:Fourier-B-t}
\end{equation}
and
\begin{equation}
\mathcal{F}\left[ B_{r}(\mathbf{r}_{0}) \right] = 
\mathcal{F}\left[ \partial_{z} \Theta(\mathbf{r}'') \right] \: 
\mathcal{F}\left[ \partial_{zz} \Psi(\mathbf{r}_{0}, \mathbf{r}'') \right] \: ,
\label{eq:Fourier-B-r}
\end{equation}
where $\mathcal{F}\left[ \cdot \right]$ denotes the 2D Fourier Transform (2DFT)
with respect to the horizontal coordinates $x$ and $y$ of the observation points.
The 2DFT's are two-dimensional functions of the wavenumbers $k_{x}$ and $k_{y}$ 
along the $x$ and $y$ directions, respectively. The wavenumbers have units of 
inverse distance. The dependence of the 2DFT's on $k_{x}$ and $k_{y}$ is conveniently
omitted to simplify the notation.

To define the 2DFT $\mathcal{F}\left[ \partial_{z} \Theta(\mathbf{r}'') \right]$ 
on the right side of equations \ref{eq:Fourier-B-t} and \ref{eq:Fourier-B-r}, 
where $\mathbf{r}'' = \left[x'' \:\: y'' \:\: z_{c} \right]^{\top}$ is the 
position vector defining a point $(x'', y'', z_{c})$ on the equivalent layer, 
we can first compute the derivative of $\Theta(\mathbf{r})$ (equation
\ref{eq:scalar-potential-Theta}) with respect to $z$ and then use 
\citet[][ p. 297]{blakely1996} to obtain:
\begin{equation}
\mathcal{F}\left[ \partial_{z} \Theta(\mathbf{r}'') \right] 
= 2\pi \, c_{m} \, e^{|k| \, z_{c}} 
\int\limits_{z_{c}}^{\infty} 
\mathcal{F}\left[ h(\mathbf{r}') \right] e^{-|k| \, z'} \:
dz' \: , \:\: z_{c} < z' \: ,
\label{eq:Fourier-Dz-Theta}
\end{equation}
where $|k| = \sqrt{k_{x}^{2} + k_{y}^{2}}$ and 
$\mathcal{F}\left[ h(\mathbf{r}') \right]$ represents the 
2DFT of the magnetization intensity distribution $h(\mathbf{r}')$
(equation \ref{eq:total-magnetization-vector}) on one horizontal slice
through the magnetic sources at $z = z'$.
In equation \ref{eq:Fourier-Dz-Theta}, we assume that 
$h(\mathbf{r}')$ is zero outside the volume $v$ of the sources
and, consequently, at any coordinate $z < z_{0} < z_{c}$.
According to \citet[][ p. 278]{blakely1996}, the 2DFT's 
$\mathcal{F}\left[ \partial_{\alpha h} \Psi(\mathbf{r}_{0}, \mathbf{r}'') \right]$ 
(equation \ref{eq:Fourier-B-alpha}),
$\mathcal{F}\left[ \partial_{th} \Psi(\mathbf{r}_{0}, \mathbf{r}'') \right]$ 
(equation \ref{eq:Fourier-B-t}) and 
$\mathcal{F}\left[ \partial_{zz} \Psi(\mathbf{r}_{0}, \mathbf{r}'') \right]$ 
(equation \ref{eq:Fourier-B-r}) are given by:
\begin{equation}
\mathcal{F}\left[ \partial_{\alpha h} \Psi(\mathbf{r}_{0}, \mathbf{r}'') \right] = 
\mathcal{U}_{\alpha} \: \mathcal{U}_{h} \: |k| \, e^{|k| \, (z_{0} - z_{c})}
\: , \:\: z_{0} < z_{c} \: , \:\: \alpha = x, y, z \: ,
\label{eq:Fourier-Dalphah-Psi}
\end{equation}
\begin{equation}
\mathcal{F}\left[ \partial_{th} \Psi(\mathbf{r}_{0}, \mathbf{r}'') \right] = 
\mathcal{U}_{t} \: \mathcal{U}_{h} \: |k| \, e^{|k| \, (z_{0} - z_{c})}
\: , \:\: z_{0} < z_{c} \: ,
\label{eq:Fourier-Dth-Psi}
\end{equation}
and
\begin{equation}
\mathcal{F}\left[ \partial_{zz} \Psi(\mathbf{r}_{0}, \mathbf{r}'') \right] 
= |k| \, e^{|k| \, (z_{0} - z_{c})} \: ,
\:\: z_{0} < z_{c} \: ,
\label{eq:Fourier-Dzz-Psi}
\end{equation}
where 
$\mathcal{U}_{x} \equiv \mathcal{U}(0^{\circ}, 0^{\circ})$,
$\mathcal{U}_{y} \equiv \mathcal{U}(0^{\circ}, 90^{\circ})$,
$\mathcal{U}_{z} \equiv \mathcal{U}(90^{\circ}, 0^{\circ})$,
$\mathcal{U}_{t} \equiv \mathcal{U}(I_{t}, D_{t})$ and
$\mathcal{U}_{h} \equiv \mathcal{U}(I_{h}, D_{h})$ are defined by the complex function
\begin{equation}
\mathcal{U}(I, D) = \hat{u}_{z}(I, D) + 
\mathrm{i} \: \frac{ \hat{u}_{x}(I, D) \: k_{x} + 
\hat{u}_{y}(I, D) \: k_{y} }{|k|} \: ,
\label{eq:U-complex-function}
\end{equation}
in terms of the imaginary unity $\mathrm{i} = \sqrt{-1}$
and the Cartesian components of the unit vector
$\hat{\mathbf{u}}(I, D)$ (equation \ref{eq:u-unit-vector}) computed with 
the corresponding inclination $I$ and declination $D$ values.
Now, by substituting the 2DFT's 
$\mathcal{F}\left[ \partial_{z} \Theta(\mathbf{r}'') \right]$ 
(equation \ref{eq:Fourier-Dz-Theta}),
$\mathcal{F}\left[ \partial_{\alpha h} \Psi(\mathbf{r}_{0}, \mathbf{r}'') \right]$
(equation \ref{eq:Fourier-Dalphah-Psi}) and 
$\mathcal{F}\left[ \partial_{th} \Psi(\mathbf{r}_{0}, \mathbf{r}'') \right]$
(equation \ref{eq:Fourier-Dth-Psi}) into those defining
$\mathcal{F}\left[ B_{t}(\mathbf{r}_{0}) \right]$
(equation \ref{eq:Fourier-B-alpha}) and
$\mathcal{F}\left[ B_{t}(\mathbf{r}_{0}) \right]$
(equation \ref{eq:Fourier-B-t}), we obtain:
\begin{equation}
\mathcal{F}\left[ B_{\alpha}(\mathbf{r}_{0}) \right] = 
2\pi \, c_{m} \, 
\mathcal{U}_{\alpha} \: \mathcal{U}_{h} \: |k| \, e^{|k| \, z_{0}} 
\int\limits_{z_{0}}^{\infty} 
\mathcal{F}\left[ h(\mathbf{r}') \right] e^{-|k| \, z'} \:
dz' \: , \:\: z_{0} < z' \: , \:\: \alpha = x, y, z \: ,
\label{eq:Fourier-B-alpha-final}
\end{equation}
and
\begin{equation}
\mathcal{F}\left[ B_{t}(\mathbf{r}_{0}) \right] = 
2\pi \, c_{m} \, 
\mathcal{U}_{t} \: \mathcal{U}_{h} \: |k| \, e^{|k| \, z_{0}} 
\int\limits_{z_{0}}^{\infty} 
\mathcal{F}\left[ h(\mathbf{r}') \right] e^{-|k| \, z'} \:
dz' \: , \:\: z_{0} < z' \: .
\label{eq:Fourier-B-t-final}
\end{equation}
Note that, in deducing equations \ref{eq:Fourier-B-alpha-final} and
\ref{eq:Fourier-B-t-final}, we have replaced $z_{c}$ by
$z_{0}$ in the 2DFT $\mathcal{F}\left[ \partial_{z} \Theta(\mathbf{r}'') \right]$ 
(equation \ref{eq:Fourier-Dz-Theta}). This change does not affect the result
because we assume that $h(\mathbf{r}')$ is zero outside the volume $v$ of the sources
and, consequently, at any coordinate $z < z_{0} < z_{c}$.

Regardless of notation differences, the 2DFT's
$\mathcal{F}\left[ B_{\alpha}(\mathbf{r}_{0}) \right]$ 
(equation \ref{eq:Fourier-B-alpha-final}) and 
$\mathcal{F}\left[ B_{t}(\mathbf{r}_{0}) \right]$ (equation \ref{eq:Fourier-B-t-final})
are exactly the same found, for example, in \citet{gunn1975} or 
\citet[][ p. 298]{blakely1996}, and are equivalent to the 2DFT's of the volume integrals 
for $B_{\alpha}(\mathbf{r})$ (equation \ref{eq:B-alpha}) and 
$B_{t}(\mathbf{r})$ (equation \ref{eq:B-t}) at the horizontal plane $z = z_{0}$. 
It can also be verified from equations 
\ref{eq:Fourier-Dz-Theta}, \ref{eq:Fourier-Dzz-Psi} and 
\ref{eq:Fourier-B-t-final} that $\mathcal{F}\left[ B_{r}(\mathbf{r}_{0}) \right]$ 
(equation \ref{eq:Fourier-B-r}) can be rewritten as follows:
\begin{equation}
\mathcal{F}\left[ B_{r}(\mathbf{r}_{0}) \right]
= \frac{1}{\mathcal{U}_{t} \,\mathcal{U}_{h}} \, 
\mathcal{F}\left[ B_{t}(\mathbf{r}_{0}) \right] \: ,
\label{eq:Fourier-B-r-final}
\end{equation}
which represents the common expression defining the RTP in Fourier domain
\cite[e.g.,][ p. 331]{blakely1996}.

Equations \ref{eq:Fourier-B-t-final} and \ref{eq:Fourier-B-r-final} establish,
for the first time, the theoretical connection between the RTP in space and Fourier domains.
We stress that these equations are valid for the particular case in which the
equivalent layer of dipoles has the same uniform magnetization direction
$\hat{\mathbf{u}}(I_{h}, D_{h})$ of the true sources 
(equation \ref{eq:total-magnetization-vector}). 


\subsection{Continuous analytical solutions with arbitrary uniform magnetization direction}
\label{subsec:arbitrary_mag}


Alternative analytical expressions for the 2DFT's 
$\mathcal{F}\left[ B_{\alpha}(\mathbf{r}_{0}) \right]$ 
(equations \ref{eq:Fourier-B-alpha} and \ref{eq:Fourier-B-alpha-final})
and
$\mathcal{F}\left[ B_{t}(\mathbf{r}_{0}) \right]$ 
(equations \ref{eq:Fourier-B-t} and \ref{eq:Fourier-B-t-final}) 
can also be deduced for the case in which the
equivalent layer has an arbitrary uniform magnetization direction 
$\hat{\mathbf{u}}(I_{p}, D_{p})$. These alternative expressions are obtained from the
simple manipulation of equations \ref{eq:Fourier-Dz-Theta},
\ref{eq:Fourier-Dalphah-Psi} and \ref{eq:Fourier-Dth-Psi}, as follows:
\begin{equation}
\mathcal{F}\left[ B_{\alpha}(\mathbf{r}_{0}) \right] = 
\mathcal{F}\left[ p(\mathbf{r}'') \right] \: 
\mathcal{F}\left[ \partial_{\alpha p} \Psi(\mathbf{r}_{0}, \mathbf{r}'') \right] \: ,
\:\: \alpha = x, y, z \: ,
\label{eq:Fourier-B-alpha-arbitrary}
\end{equation}
and
\begin{equation}
\mathcal{F}\left[ B_{t}(\mathbf{r}_{0}) \right] = 
\mathcal{F}\left[ p(\mathbf{r}'') \right] \: 
\mathcal{F}\left[ \partial_{tp} \Psi(\mathbf{r}_{0}, \mathbf{r}'') \right] \: ,
\label{eq:Fourier-B-t-arbitrary}
\end{equation}
where
\begin{equation}
\mathcal{F}\left[ p(\mathbf{r}'') \right] 
= \frac{\mathcal{U}_{h}}{\mathcal{U}_{p}} \, 
\mathcal{F}\left[ \partial_{z} \Theta(\mathbf{r}'') \right] \: ,
\label{eq:Fourier-p}
\end{equation}
\begin{equation}
\mathcal{F}\left[ \partial_{\alpha p} \Psi(\mathbf{r}_{0}, \mathbf{r}'') \right] = 
\mathcal{U}_{\alpha} \: \mathcal{U}_{p} \: |k| \, e^{|k| \, (z_{0} - z_{c})}
\: , \:\: z_{0} < z_{c} \: , \:\: \alpha = x, y, z \: ,
\label{eq:Fourier-Dalphap-Psi}
\end{equation}
and
\begin{equation}
\mathcal{F}\left[ \partial_{tp} \Psi(\mathbf{r}_{0}, \mathbf{r}'') \right] = 
\mathcal{U}_{t} \: \mathcal{U}_{p} \: |k| \, e^{|k| \, (z_{0} - z_{c})}
\: , \:\: z_{0} < z_{c} \: ,
\label{eq:Fourier-Dtp-Psi}
\end{equation}
with $\mathcal{U}_{p} \equiv \mathcal{U}(I_{p}, D_{p})$ (equation 
\ref{eq:U-complex-function}) defined in terms of the arbitrary constant 
inclination $I_{p}$ and declination $D_{p}$ values.
Note that, in this case, it is not possible to compute the 2DFT 
$\mathcal{F}\left[ B_{r}(\mathbf{r}_{0}) \right]$ of the RTP anomaly from the
2DFT $\mathcal{F}\left[ p(\mathbf{r}'') \right]$ of the generic physical-property 
distribution $p(\mathbf{r}'')$ because, according to equations \ref{eq:Fourier-B-r}, 
\ref{eq:Fourier-Dzz-Psi} and \ref{eq:Fourier-p}:
\begin{equation}
\mathcal{F}\left[ B_{r}(\mathbf{r}_{0}) \right]
\ne \mathcal{F}\left[ p(\mathbf{r}'') \right] \, 
\mathcal{F}\left[ \partial_{zz} \Psi(\mathbf{r}_{0}, \mathbf{r}'') \right] \: .
\label{eq:Fourier-B-r-wrong}
\end{equation}

The alternative analytical expressions presented here 
(equations \ref{eq:Fourier-B-alpha-arbitrary}--\ref{eq:Fourier-B-r-wrong})
show, in Fourier domain, that the planar equivalent layer with 
generic magnetic moment distribution 
$p(\mathbf{r}'')$ obtained by using an arbitrary uniform magnetization direction
$\hat{\mathbf{u}}(I_{p}, D_{p})$ can exactly reproduce the magnetic field components
$B_{\alpha}(\mathbf{r}_{0})$ and approximated total-field anomaly
$B_{t}(\mathbf{r}_{0})$, but cannot be used to compute the RTP anomaly
$B_{r}(\mathbf{r}_{0})$.


\subsection{Analytical solutions in discrete form}
\label{subsec:discrete_forms}

$\mathbf{r}_{i} = \left[x_{i} \:\: y_{i} \:\: z_{i} \right]^{\top}$

$\mathbf{r}''_{j} = \left[x''_{j} \:\: y''_{j} \:\: z_{c} \right]^{\top}$

\begin{equation}
V(\mathbf{r}_{i}) \approx - \sum_{j=1}^{M} p_{j} \:\: \partial_{h} \Psi(\mathbf{r}_{i}, \mathbf{r}''_{j}) \: ,
\label{eq:V-Neumann-i}
\end{equation}

\begin{equation}
B_{t}(\mathbf{r}_{i}) \approx \sum_{j=1}^{M} p_{j} \:\: \partial_{tp} \Psi(\mathbf{r}_{i}, \mathbf{r}''_{j}) \: ,
\label{eq:B-t-Neumann-i}
\end{equation}

\begin{equation}
B_{r}(\mathbf{r}_{i}) \approx \sum_{j=1}^{M} p_{j} \:\: \partial_{zz} \Psi(\mathbf{r}_{i}, \mathbf{r}''_{j}) \: ,
\label{eq:B-r-Neumann-i}
\end{equation}

\begin{equation}
p_{j} = \partial_{z} \Theta(\mathbf{r}''_{j}) \: \Delta s_{j}
\label{eq:pj}
\end{equation}

\begin{equation}
\mathbf{d} \approx \mathbf{A} \: \mathbf{p}
\label{eq:discrete-space-domain}
\end{equation}



1 - apresentar versões discretas das integrais \ref{eq:V-Neumann}, \ref{eq:B-Neumann},
\ref{eq:B-t-Neumann} e \ref{eq:B-r-Neumann}

2 - usar a formulação BTTB para definir estas versões discretas no domínio de Fourier

3 - estabelecer que as expressões BTTB são as versões numéricas das equações 
\ref{eq:Fourier-B-t} e \ref{eq:Fourier-B-r}. Mencionar que estas versões numéricas 
no domínio de Fourier são ligeiramente diferentes das expressões analíticas porque
estas últimas não incluem os efeitos de discretização e truncamento.

\section{Numerical simulations}
\label{sec:numerical_simulations}

1 - gerar dz theta em um plano

2 - usar o dz theta discreto para validar as eqs definidas no item 1 da sec anterior

3 - calcular as 2DFT's discretas para validar as eqs definidas no ítem 2 da sec anterior


%\section{Sources with nonuniform magnetization direction}
%\label{sec:eqlayer_nonuniform}
%
%
%By combining equations \ref{eq:V-summation} and \ref{eq:V-l-Neumann}, we obtain
%that the resultant magnetic scalar potential $V(\mathbf{r})$ 
%(equation \ref{eq:V-volume-integral}) can be represented by
%the superposition of equivalent layers associated with all magnetization components:
%\begin{equation}
%V(\mathbf{r}) = -
%\int\limits_{-\infty}^{\infty} \int\limits_{-\infty}^{\infty}
%\frac{\| \mathbf{v}(\mathbf{r}'') \|}{2\pi} \:
%\nabla \frac{1}{\| \mathbf{r} - \mathbf{r}'' \|}^{\top}
%\hat{\mathbf{v}}(\mathbf{r}'') \: 
%dx''dy'' \: , \quad z < z_{c} \: ,
%\label{eq:V-superposition-layers}
%\end{equation}
%where
%\begin{equation}
%\mathbf{v}(\mathbf{r}'') = 
%\begin{bmatrix}
%v_{x}(\mathbf{r}'') \\ v_{y}(\mathbf{r}'') \\ v_{z}(\mathbf{r}'')
%\end{bmatrix} =
%\sum\limits_{\ell = 1}^{L}
%\partial_{z} \Theta_{\ell}(\mathbf{r}'') \:
%\hat{\mathbf{u}}(I_{\ell}, D_{\ell}) \quad ,
%\label{eq:v-vector}
%\end{equation}
%and
%\begin{equation}
%\hat{\mathbf{v}}(\mathbf{r}'') = 
%\begin{bmatrix}
%\hat{v}_{x}(\mathbf{r}'') \\ 
%\hat{v}_{y}(\mathbf{r}'') \\ 
%\hat{v}_{z}(\mathbf{r}'')
%\end{bmatrix} =
%\frac{\mathbf{v}(\mathbf{r}'')}{\| \mathbf{v}(\mathbf{r}'') \|} \quad .
%\label{eq:v-unit-vector}
%\end{equation}
%Similarly, equations \ref{eq:B-summation}, \ref{eq:Delta-T}, \ref{eq:B-l-Neumann} and
%\ref{eq:Delta-T-l-Neumann} can be used to show that the 
%resultant magnetic induction field $\mathbf{B}(\mathbf{r})$ and
%approximated total-field anomaly $\Delta\tilde{T}(\mathbf{r})$ also have the 
%following equivalent-layer representations:
%\begin{equation}
%\mathbf{B}(\mathbf{r}) = 
%\int\limits_{-\infty}^{\infty} \int\limits_{-\infty}^{\infty}
%\frac{\| \mathbf{v}(\mathbf{r}'') \|}{2\pi} \: 
%\mathbf{G}(\mathbf{r}, \mathbf{r}'') \: \hat{\mathbf{v}}(\mathbf{r}'') \: 
%dx''dy'' \: , \:\: z < z_{c} \: ,
%\label{eq:B-superposition-layers}
%\end{equation}
%and
%\begin{equation}
%\Delta\tilde{T}(\mathbf{r}) = 
%\int\limits_{-\infty}^{\infty} \int\limits_{-\infty}^{\infty}
%\frac{\| \mathbf{v}(\mathbf{r}'') \|}{2\pi} \: 
%\hat{\mathbf{u}}(I_{t}, D_{t})^{\top} 
%\mathbf{G}(\mathbf{r}, \mathbf{r}'') \: \hat{\mathbf{v}}(\mathbf{r}'') \: 
%dx''dy'' \: , \:\: z < z_{c} \: .
%\label{eq:Delta-T-superposition-layers}
%\end{equation}
%
%Note that the equivalent layer representing the resultant magnetic scalar potential 
%$V(\mathbf{r})$ (equation \ref{eq:V-superposition-layers}),
%magnetic induction field $\mathbf{B}(\mathbf{r})$ (equation \ref{eq:B-superposition-layers})
%and approximated total-field anomaly $\Delta\tilde{T}(\mathbf{r})$
%(equation \ref{eq:Delta-T-superposition-layers}) does not have a uniform
%magnetization direction. Rather it has a spatially varying direction 
%$\hat{\mathbf{v}}(\mathbf{r}'')$ (equations \ref{eq:v-vector} and \ref{eq:v-unit-vector})
%that combines the individual uniform directions $\hat{\mathbf{u}}(I_{\ell}, D_{\ell})$
%of all $L$ magnetization components. 
%Besides, this layer has a strictly positive magnetic-moment distribution 
%$\tfrac{\| \mathbf{v}(\mathbf{r}'') \|}{2\pi}$ defined as the ratio of the
%Euclidean norm of $\mathbf{v}(\mathbf{r}'')$ (equation \ref{eq:v-vector}) and the positive
%constant $2\pi$. 
%Equations \ref{eq:V-superposition-layers}--\ref{eq:Delta-T-superposition-layers} generalize, 
%in the space domain, the all-positive 
%magnetic-moment distribution presented by \cite{reis_etal2020} to the case in which the sources 
%are arbitrarily magnetized with heterogeneous magnetization intensity and direction.
%
%
%\subsection{Relationships between the equivalent layer and the true sources}
%\label{subsec:relationships_true_sources}
%
%
%Note that $\Theta_{\ell}(\mathbf{r})$ (equation \ref{eq:Theta-l}) 
%and $\partial_{z} \Theta_{\ell}(\mathbf{r}'')$ 
%(equations \ref{eq:Theta-l-Neumann} and \ref{eq:V-l-Neumann} and \ref{eq:v-vector}) 
%are analogous to the gravitational potential and 
%vertical component of the gravitational attraction (gravity disturbance) that 
%would be produced by the $\ell$th magnetization component at the observation point
%$\mathbf{r}$ and at the point $\mathbf{r}''$ on the equivalent layer plane $z_{c}$,
%respectively. 
%It is also important noting that these two functions are strictly positive because
%they are defined in terms of the positive total-magnetization intensity distribution 
%$h_{\ell}(\mathbf{r}')$ (equation \ref{eq:h-vector-total}) of the $\ell$th magnetization 
%component.